% lecture07.tex - 叉积及应用
% 第 7 章: 叉积及应用 (Cross Products and Applications)

\chapter{叉积及应用}
\label{ch:crossproduct}

\begin{abstract}
叉积 (cross product, 又称向量积 vector product) 是一种运算, 它由三维空间中的两个向量生成一个同时垂直于这两个向量的向量。在本章中, 我们将学习它的定义、性质, 以及它与行列式之间的深刻联系。
\end{abstract}

% ============================================
\section{二维空间中的叉积}
% ============================================

\begin{definition}[二维叉积]
对于 $\Rtwo$ 中的两个向量:
\[
\vv \times \vw = v_1 w_2 - v_2 w_1 = \det\twomat{v_1}{w_1}{v_2}{w_2}
\]
结果是一个\textbf{数 (number)}, 而不是向量。
\end{definition}

\begin{intuition}
二维叉积 (2D cross product) = 由两个向量张成的平行四边形的有向面积

正号: $\vw$ 在 $\vv$ 的左侧\\
负号: $\vw$ 在 $\vv$ 的右侧
\end{intuition}

\begin{center}
\begin{tikzpicture}[scale=1.2]
    % 平行四边形
    \fill[blue!20] (0,0) -- (2,0.5) -- (3,2) -- (1,1.5) -- cycle;
    \draw[vector, red!70, very thick] (0,0) -- (2,0.5) node[midway, below] {$\vv$};
    \draw[vector, green!60!black, very thick] (0,0) -- (1,1.5) node[midway, left] {$\vw$};
    \draw[dashed] (2,0.5) -- (3,2);
    \draw[dashed] (1,1.5) -- (3,2);
    \node at (1.5,1) {面积 $= |\vv \times \vw|$};
\end{tikzpicture}
\end{center}

% ============================================
\section{三维空间中的叉积}
% ============================================

\begin{definition}[三维叉积]
对于 $\vv = \threevec{v_1}{v_2}{v_3}$ 和 $\vw = \threevec{w_1}{w_2}{w_3}$:
\[
\vv \times \vw = \threevec{v_2 w_3 - v_3 w_2}{v_3 w_1 - v_1 w_3}{v_1 w_2 - v_2 w_1}
\]
\end{definition}

\begin{theorem}[行列式公式]
\[
\vv \times \vw = \det\threemat{\vi}{\vj}{\vk}{v_1}{v_2}{v_3}{w_1}{w_2}{w_3}
\]
(按第一行作形式展开)
\end{theorem}

\begin{intuition}
叉积 $\vv \times \vw$:
\begin{itemize}
    \item \textbf{方向 (direction):} 同时垂直于两个向量 (右手定则 right-hand rule)
    \item \textbf{大小 (magnitude):} 由 $\vv$ 和 $\vw$ 张成的平行四边形的面积
\end{itemize}
\end{intuition}

\begin{center}
\begin{tikzpicture}[scale=1]
    % 三维坐标示意
    \draw[axis] (0,0) -- (3,0) node[right] {$x$};
    \draw[axis] (0,0) -- (0,3) node[above] {$z$};
    \draw[axis] (0,0) -- (-1,-0.7) node[below left] {$y$};

    % 向量
    \draw[vector, red!70, very thick] (0,0) -- (2,-0.5) node[right] {$\vv$};
    \draw[vector, green!60!black, very thick] (0,0) -- (-0.5,0) node[left] {$\vw$};
    \draw[vector, blue!70, ultra thick] (0,0) -- (0,2) node[above] {$\vv \times \vw$};

    % 平行四边形底面
    \fill[gray!20] (0,0) -- (2,-0.5) -- (1.5,-0.5) -- (-0.5,0) -- cycle;
\end{tikzpicture}
\end{center}

% ============================================
\section{叉积的性质}
% ============================================

\begin{theorem}[主要性质]
\begin{enumerate}
    \item \textbf{反交换律 (anti-commutativity):} $\vv \times \vw = -(\vw \times \vv)$
    \item \textbf{分配律 (distributivity):} $\vv \times (\vw + \vu) = \vv \times \vw + \vv \times \vu$
    \item \textbf{标量乘法 (scalar multiplication):} $(c\vv) \times \vw = c(\vv \times \vw)$
    \item \textbf{正交性 (orthogonality):} $\vv \times \vw \perp \vv$ 且 $\vv \times \vw \perp \vw$
    \item \textbf{自身叉积为零 (self-cross is zero):} $\vv \times \vv = \vzero$
\end{enumerate}
\end{theorem}

\begin{warning}
叉积\textbf{不满足交换律}!
\[
\vv \times \vw \neq \vw \times \vv
\]
事实上, 它们互为相反向量。
\end{warning}

\begin{theorem}[模长公式]
\[
\norm{\vv \times \vw} = \norm{\vv} \norm{\vw} \sin\theta
\]
其中 $\theta$ 是两个向量之间的夹角。
\end{theorem}

% ============================================
\section{从线性变换视角看叉积}
% ============================================

\begin{intuition}
更深刻的看法: 叉积 (cross product) 可以通过对偶 (duality) 来定义。

函数 $f(\vx) = \det[\vv \,|\, \vw \,|\, \vx]$ 是关于 $\vx$ 的线性函数 (linear function)。由对偶性, 必存在一个向量 $\vp$ 使得:
\[
f(\vx) = \vp \cdot \vx
\]
这个 $\vp$ 正是 $\vv \times \vw$!
\end{intuition}

\begin{theorem}
\[
\det[\vv \,|\, \vw \,|\, \vx] = (\vv \times \vw) \cdot \vx
\]
\end{theorem}

% ============================================
\section{基向量的叉积}
% ============================================

\begin{theorem}[基向量的叉积]
\begin{align*}
\vi \times \vj &= \vk & \vj \times \vi &= -\vk \\
\vj \times \vk &= \vi & \vk \times \vj &= -\vi \\
\vk \times \vi &= \vj & \vi \times \vk &= -\vj
\end{align*}
\end{theorem}

\begin{example}
\[
\threevec{2}{3}{4} \times \threevec{5}{6}{7}
\]
\begin{align*}
&= (3 \times 7 - 4 \times 6)\vi - (2 \times 7 - 4 \times 5)\vj + (2 \times 6 - 3 \times 5)\vk \\
&= (21-24)\vi - (14-20)\vj + (12-15)\vk \\
&= -3\vi + 6\vj - 3\vk = \threevec{-3}{6}{-3}
\end{align*}
\end{example}

% ============================================
\section{应用}
% ============================================

\begin{practical}
\textbf{物理: 力矩 (torque)}

力 $\vec{F}$ 关于距离为 $\vec{r}$ 的一点的力矩为:
\[
\vec{\tau} = \vec{r} \times \vec{F}
\]
力矩的方向垂直于力与力臂所在的平面。
\end{practical}

\begin{practical}
\textbf{物理: 洛伦兹力 (Lorentz force)}

在磁场中运动的带电粒子所受的力为:
\[
\vec{F} = q\vec{v} \times \vec{B}
\]
\end{practical}

\begin{practical}
\textbf{计算机图形学: 曲面法向量 (surface normal)}

求顶点为 $A$, $B$, $C$ 的三角形面的法向量 (normal vector):
\[
\vec{n} = (\vec{B} - \vec{A}) \times (\vec{C} - \vec{A})
\]
\end{practical}

\begin{practical}
\textbf{三角形面积计算}

顶点为 $A$, $B$, $C$ 的三角形的面积为:
\[
\text{Area} = \frac{1}{2}\norm{(\vec{B}-\vec{A}) \times (\vec{C}-\vec{A})}
\]
\end{practical}

% ============================================
\section{三重积}
% ============================================

\begin{definition}[标量三重积 (scalar triple product)]
\[
\va \cdot (\vb \times \vc) = \det[\va \,|\, \vb \,|\, \vc]
\]
结果 = 由三个向量张成的平行六面体的有向体积
\end{definition}

\begin{theorem}[轮换性质]
\[
\va \cdot (\vb \times \vc) = \vb \cdot (\vc \times \va) = \vc \cdot (\va \times \vb)
\]
\end{theorem}

% ============================================
\section{练习}
% ============================================

\begin{exercise}
计算叉积:
\[
\va = \threevec{1}{2}{3}, \quad \vb = \threevec{4}{5}{6}
\]
\end{exercise}

\begin{exercise}
求以 $\va = \threevec{1}{0}{0}$ 和 $\vb = \threevec{1}{1}{0}$ 为边的平行四边形的面积。
\end{exercise}

\begin{exercise}
证明 $\va \times \vb$ 同时垂直于 $\va$ 和 $\vb$。
\end{exercise}

\begin{exercise}
求过点 $(1,0,0)$, $(0,1,0)$, $(0,0,1)$ 的平面的法向量。
\end{exercise}

\begin{exercise}[挑战]
证明:
\[
\norm{\va \times \vb}^2 + (\va \cdot \vb)^2 = \norm{\va}^2 \norm{\vb}^2
\]
(这称为拉格朗日恒等式 Lagrange's identity)
\end{exercise}

\begin{problem}
计算以以下向量为棱的平行六面体的体积:
\[
\va = \threevec{1}{1}{0}, \quad \vb = \threevec{0}{1}{1}, \quad \vc = \threevec{1}{0}{1}
\]
\end{problem}
